\chapter{68}




«Wie ein waschechter Schweizer siehst du nicht gerade aus,» Avas Lachen erreichte ihn 
über den Tisch, «Ciro, spanischer Vorname, hm, lass mich raten.»

Sie kniff nachdenklich die Augen zusammen. 
Louis setzte seine Sonnenbrille auf, sah in ihr Gesicht und blieb still.
Mit ihren Mutmassungen um seine Identität landete sie schliesslich
beinahe einen Treffer.

«Dann also, Schweizer Vater, Mutter Spanierin?»

Den Kopf vorgestreckt, wartete sie neugierig. Louis erlöste sie.

«Nun tauschst du die Nationalitäten aus und dann stimmt’s.»

Ava klatschte in die Hände.

Die weiteren Gespräche drehten sich um die Schafe. Um die Arbeiten, die
anstanden, hier oben und unten auf der \emph{Sunnewaid}. Um die
bevorstehende \emph{Schafscheid} und die anschliessende \emph{Schur}.
Diese stand offensichtlich unter Manuelas Regie. Joh hatte sich
zurückgelehnt und ihr konzentriert zugehört. 
Die ganze Zeit über suchte Ava Louis’ Blick. Wiederholt einen Moment zu
lange. Was er gewöhnlich als schmeichelhaft empfand und genoss, war ihm
unter den gegebenen Umständen unangenehm. Nicht zuletzt auch, weil er
glaubte, jedes Mal einen undefinierbaren Schatten über Avas Gesicht huschen
zu sehen.

fiSelten hatte er sich im Umgang mit neuen Menschen so unsicher gefühlt.
Er sehnte sich nach Oberflächlichkeit. Nach banalem Austausch von Worten und ohne
Erwartung an Intimität. Ciro war nichts anderes als eine Spielfigur und er rutschte immer
wieder aus der Rolle. In heiklen Situationen wusste er nicht, ob ein
passendes Verhalten näher bei Louis oder näher bei Ciro lag. Es war anstrengend und schon fast unheimlich.

\sectionbreak

Den Abstieg hatten sie gut gemeistert, Mazi und er. Sie hatten die letzte Gondel erwischt. Abends sassen sie am grossen Tisch vor der \emph{Sunnewaid}. Sie sprachen kaum. Es passte.
Louis war, als geniesse auch Mazi die Stille.

\chapter{69} 

\section{Bettmeralp (CH), Samstag, 1.
September}

Louis’ erster Blick in den nächsten Morgen war verblüffend angenehm. Er
hatte geschlafen wie ein Baby. 
Der neue Tag würde gut gefüllt sein. Es war ihm recht. Er kannte die Aufgaben, 
die leicht zu erledigen waren. Er würde immer wieder seinen Gedanken nachhängen können und  
war endlich wieder bei seiner \emph{Gibson}. Wann immer möglich, würde er sich Pausen gönnen 
und sich mit ihr zurückziehen. Manuela kam erst in zwei Tagen zurück auf den Hof.

Er hängte die nasse Wäsche an die Leinen unter die Sonne. Nach dem
Setzen der Klammern hob er den Wäschekorb hoch und hielt inne. Hingerissen schaute er abwechselnd links und rechts in das Weiss der Bettbezüge und zwischen den vollgehängten Leinen entlang bis nach vorne.
Ihm war, als stehe er mitten in einem Duftkorridor. Einem betörenden
Cocktail aus Gras, Wiesenblumen, Waschmittel und warmer Haut, feucht und
kühl. Und mit der Behaglichkeit kam die Überzeugung. Er musste in
Bewegung kommen, einen Schritt machen, egal welchen, bevor die Angst zurückkam.
Bevor er noch tiefer in diesen unheilvollen Gefühlssumpf geriet, der ihn mit jeder seiner Handlungen weiter nach unten ziehen und verschlingen  würde. Er entschied, Valentina anzurufen.
